% ============================================================
%  第二章
% ============================================================
\section{从 $U(1)$ 到非交换：磁荷究竟属于谁}
\markboth{第二章\ 从 $U(1)$ 到非交换：磁荷究竟属于谁}{第二章\ 从 $U(1)$ 到非交换：磁荷究竟属于谁}

\noindent{\small\color{shuyiGray}对应底本：原 p.7--p.12。}\par\medskip

\srcp{7}让我先解释一点 Goddard、Nuyts 与 Olive 说过的东西。我们考虑规范理论，先从阿贝尔规范群 $G=U(1)$ 的情形开始——在现实世界里，这就是电磁作用的规范群。有像电子这样的“带电粒子”。一个粒子的电电荷，最根本的定义是：它是与这个粒子相伴的 $G$ 的一个表示；不过，由于 $U(1)$ 的每个表示都是某个基本表示的 $n$ 次幂（$n$ 为某个整数），你可以把电电荷就想成是一个整数。一种不那么抽象的说法是：电电荷是“量子化”的，任何粒子的电荷都是电子电荷的整数倍。

\subsection{狄拉克单极子}

\srcp{8}在 $U(1)$ 规范理论里，“规范场”是时空上一个酉复线丛 $\cL$ 上的联络 $A$。记 $F=\dd A$ 为 $A$ 的曲率。电子的量子波函数是 $\cL$ 的一个截面 $\psi$。狄拉克在 1930 年代引入了磁单极子的想法（作为量子对象；经典物理里的想法还要更早）。狄拉克的单极子概念是：单极子的位置是空间中的一个奇点 $p$，在 $p$ 处线丛 $\cL$ 与联络 $A$ 没有定义；$\RR^3$ 去掉点 $p$ 之后可缩到一个二维球面 $S$，而磁荷是
\[
  m=\int_S c_1(\cL),
\]
是一个整数。

\begin{figure}[htbp]
\centering
\begin{tikzpicture}[scale=1.15]
  \shade[ball color=gray!6] (0,0) circle (1.1);
  \draw[thick] (0,0) circle (1.1);
  \draw[densely dashed] (-1.1,0) arc (180:360:1.1 and 0.38);
  \draw (-1.1,0) arc (180:0:1.1 and 0.38);
  \fill (0.22,0.3) circle (1.7pt);
  \node[right=3pt] at (0.22,0.3) {\small $p$};
  \node[above] at (0,1.22) {\small $S$};
  \draw[->,gray] (1.0,-0.75) -- (0.55,-0.2);
  \node[right,gray] at (1.0,-0.78) {\scriptsize $\cL,\,A$ 在 $p$ 处无定义};
\end{tikzpicture}
\caption{狄拉克单极子。空间 $\RR^3$ 中挖去一点 $p$，围绕它的二维球面 $S$ 上，线丛 $\cL$ 的第一陈类给出磁荷。$\cL$ 在 $p$ 处没有定义，因此不存在整体的 $A$；所谓“狄拉克弦”只是局部截面失效的地方，并非物理的弦。}
\label{fig:dirac}
\end{figure}

\srcp{9}如果我们要求联络 $A$ 满足麦克斯韦方程，就能确定静止单极子处的联络必须是什么。曲率是
\begin{equation}
  F=\frac{m}{2}\star \dd\frac{1}{|\vec x|}.
  \label{eq:monopole}
\end{equation}
当且仅当 $m$ 是整数时，$\RR^3\setminus p$ 上存在线丛 $\cL$ 与联络 $A$，使得曲率 $F=\dd A$ 就是我写下的这个；而 $\cL$ 与 $A$ 在相差同构的意义下唯一。

\subsection{非阿贝尔理论：电是表示，磁是拓扑}

\srcp{10}现在讨论非阿贝尔规范理论，比如规范群取一个单的非阿贝尔群 $G$。时空上现在有一个 $G$-丛 $E$，带联络 $A$ 与曲率 $F=\dd A+A\wedge A$。电电荷容易定义：任何粒子的电电荷都是 $G$ 的一个表示 $R$。磁荷该是什么意思，则比较微妙。对 $G=U(1)$ 我说过的话，最直接的类比是：单极子仍是空间中的一点（或一小块区域），在那里 $E$ 与 $A$ 没有定义；而磁荷分类的是围绕它的二维球面 $S$ 上 $G$-丛的选择。按这个解释（在合适的语境里它是恰当的），$G$ 的磁荷取值于 $\pi_1(G)$。

\srcp{11}这不是 GNO 想要的答案。他们考虑的是一种理论：通过希格斯机制（就是 2012 年被 CERN 的 LHC 加速器著名地证实的那种希格斯机制），规范群在低能下、或对大多数实际目的而言，被有效地约化到极大环面 $T\subset G$。如果这种约化只在远离点 $p$ 的地方发生，那么按前面同样的推理，我们可以定义一个磁荷，它分类的是二维球面 $S$ 上 $T$-丛（而不是 $G$-丛）的选择，因此取值于 $\pi_1(T)$；而对环面 $T$ 来说，这同一个东西就是 $\Hom(U(1),T)$。

\srcp{12}另一方面，一旦假定规范对称性从 $G$ 有效地约化到极大环面 $T$，我们就应当用 $T$ 的表示来分类电电荷。阿贝尔群 $T$ 的不可约表示都是一维的，所以它们对应于 $\Hom(T,U(1))$ 的元素。于是在 GNO 考虑的框架里，他们有：
\begin{itemize}
  \item 电电荷：$\Hom(T,U(1))$；
  \item 磁荷：$\Hom(U(1),T)$。
\end{itemize}

\begin{table}[htbp]
\centering\small
\caption{电与磁：从 $U(1)$ 到非阿贝尔群，再到希格斯化后的环面。}
\label{tab:electric-magnetic}
\begin{tabular}{@{}lll@{}}
\toprule
情形 & 电电荷 & 磁荷 \\
\midrule
$G=U(1)$ & $n\in\ZZ$（电荷） & $m=\int_S c_1(\cL)\in\ZZ$ \\
非阿贝尔单群 $G$ & $G$ 的表示 $R$ & 朴素看法：$\pi_1(G)$ \\
$G$ 破缺到 $T$ & $\Hom(T,U(1))=X^*(T)$ & $\Hom(U(1),T)=X_*(T)$ \\
\bottomrule
\end{tabular}
\end{table}


\begin{yizhu}[monopole、magnetic charge]
\textit{monopole} 在物理文献里常译“磁单极子”；本书按数译术语库统一作“单极子”，
只在原文出现 \textit{magnetic monopole}、需要强调“磁”的时候写“磁单极子”。
\textit{magnetic charge} 一律译“磁荷”。公式里的记号随原文：电电荷记 $e$，磁荷在
第 2 章记 $m$，在第 8 章记 $\rho$ 所代表的余权。
\end{yizhu}

\begin{pizhu}{狄拉克量子化：磁荷为什么是整数}
设 $S$ 是包围单极子的球面。线丛 $\cL\to S$ 由转移函数 $g:S^1\cap S^1\to U(1)$ 的缠绕数 $m$ 分类，故 $c_1(\cL)=m$，而 $\int_S F=2\pi m$（这里把 $F/2\pi$ 看作第一陈类的 de Rham 代表）。式 \eqref{eq:monopole} 中的系数 $\tfrac m2$ 正是按这个归一化选的：在 $\RR^3\setminus 0$ 上，$\star\dd(1/r)$ 在单位球面上的通量恰为 $-4\pi$，所以 $F=\tfrac m2\star\dd(1/r)$ 给出 $|\int_S F|=2\pi|m|$，差一个符号取决于 $S$ 的定向。

\smallskip
\textbf{与电荷的关系。}假设有一个电荷为 $e$ 的粒子绕单极子一圈，其波函数要乘以 $\exp(\ii e\oint A)$。环绕路径 $\gamma$ 可以缩成一个以 $S$ 的一部分为边界的曲面 $\Sigma$，而 $\oint_\gamma A=\int_\Sigma F$。要求这个相位与曲面的选择无关，就得到
\[
  e\int_S F\in 2\pi\ZZ
  \quad\Longrightarrow\quad
  e\,m\in 4\pi\ZZ\quad(\text{本节的归一化}).
\]
这就是狄拉克量子化：电与磁的乘积必须落在 $\ZZ$ 的格点上。若把电子的电荷归一为 $1$，就等于把 $4\pi$ 吸收进了单位——Witten 说“电荷是整数”，说的正是这个归一化。真正与单位无关的陈述是一个配对：$X^*(T)\otimes X_*(T)\to \ZZ$。下一章会看到，这个配对正是 GNO 对偶的全部内容。
\end{pizhu}

\begin{pizhu}{$\pi_1(G)$ 为什么不够：从奇点到光滑孤子}
把 $S^2$ 分成两个圆盘，其交线是赤道 $S^1$；$G$-丛在 $S^2$ 上的同构类由转移函数 $S^1\to G$ 的同伦类给出，即 $\pi_1(G)$。对单连通的 $G$（如 $SU(N)$），$\pi_1(G)=0$，于是球面上的一切 $G$-丛都平凡：在这个意义上，“点状奇点”式的非阿贝尔单极子根本不存在。

这正是 't Hooft--Polyakov 单极子的要害：它之所以是光滑的，是因为希格斯场在无穷远处把 $G$ 破缺到 $U(1)$，无穷远处的球面上真正被分类的丛是 $U(1)$-丛，其第一陈类给出整数磁荷；而在原点附近，场是处处有定义的。换句话说，非阿贝尔单极子的磁荷不是“挖掉一个点”的拓扑，而是“破缺之后剩下的 $U(1)$”的拓扑。

对环面 $T=U(1)^n$，$\pi_1(T)=\ZZ^n=X_*(T)$，同构由 $k\mapsto(z\mapsto z^k)$ 给出；这就解释了为什么第 12 页上磁荷是 $\Hom(U(1),T)$。\emph{余特征标}（coweight）这个词的来源也在这里：它是 $U(1)$ 到 $T$ 的映射，而不是 $T$ 到 $U(1)$ 的映射。
\end{pizhu}

\begin{pizhu}{三种“同一性”}
本章已经把三种非常不同的东西摆在了同一个位置：群 $\pi_1(G)$ 的拓扑、球面上丛的分类、以及 $\Hom(U(1),T)$ 的表示论。Witten 后面所有的话，都是在这三者的“同一性”上做文章。读这一章时值得记住一个反差：$U(1)$ 的情形是纯粹的陈类，$SU(2)$ 破缺到 $U(1)$ 的情形是纯粹的狄拉克量子化，而一般的 $G$ 破缺到 $T$ 的情形——GNO 的情形——则是这两者的混合。第 3 章将说明，这个混合体的群论名字叫“根数据的对偶”。
\end{pizhu}
